\section{Impact}
\label{sec:discussion}

\subsection{What the tool enables}
\tool changes how graph-based targeted-attack detectors can be studied by supplying an open,
controllable, confounder-controlled evaluation instrument. Its impact is as \emph{research
infrastructure} rather than a single new score. Most directly, it is the instrument behind the
authors' companion studies on graph-based spear-phishing and BEC detection (in preparation),
which use it to run the analyses A1--A4---which structures carry signal, whether an advantage
is mechanism or leakage, how detection degrades under an adaptive attacker, and whether
robustness transfers across conditions---reproducibly and under audited conditions. Beyond that
programme, it improves the pursuit of \emph{existing} questions for the community: for
COMPA-style heuristics \cite{egele2013compa}, lateral-phishing features
\cite{ho2019detecting}, and GNN fraud detectors \cite{dou2020enhancing}, it offers a common,
leakage-audited harness so that claims become falsifiable and comparable rather than
optimistic point estimates on private data \cite{arp2022dos}. Third, it lowers the cost of
contribution: a new detector, attack, or graph is a few lines against one interface, so a
group can drop in its own model and obtain a leakage-controlled comparison the same day.

\subsection{Threats to validity and limitations}
The most important limitation is \emph{synthetic realism}. The procedural organization model
captures hierarchy, partners, and rhythm, but not the full messiness of real correspondence;
we mitigate this by supporting real graphs (Enron, SNAP) and by using synthetic-to-real
concordance (Section~\ref{sec:results}) as an explicit validity check, but the tool validates
\emph{mechanism and topology}, not the detection of any specific real-world incident.
Second, the Enron and SNAP graphs supply real \emph{structure} but the attack \emph{labels}
are injected, not observed; \tool is therefore a controlled testbed, not a substitute for a
labelled operational corpus, and results should be read as evidence about detector behaviour
under known conditions. Third, the detector panel, while spanning heuristics to temporal
GNNs, is not exhaustive; the single-interface design is precisely what makes extending it
inexpensive. Finally, the leak-aware controls remove the confounders we identified (rhythm,
degree, volume, identity), but no control set is complete---the tool makes leakage
\emph{auditable and explicit}, which is the prerequisite for finding the next confounder,
not a guarantee of its absence.

\subsection{Ethics and data}
All experiments use either a fully synthetic population, which contains no personal data and
raises no GDPR concerns, or public graphs whose structure is real but whose attack labels are
controlled injections. The toolkit is designed for defensive research---measuring and
improving detectors---and does not generate deployable attack content. This makes \tool
freely shareable in settings where real phishing corpora, bound by privacy and
confidentiality, cannot be distributed.